% CHAPTER 2: LITERATURE REVIEW
\chapter{Literature Review}
\label{ch:litreview}

This chapter provides a comprehensive survey of the literature relevant to
financial fraud detection using multi-paradigm deep learning. We begin by
establishing the domain background, including the nature of financial fraud, the
characteristics of benchmark datasets, and the fairness considerations that
motivate responsible model design. We then review the state-of-the-art in
software solutions and frameworks that support fraud detection research. The
core of the chapter presents a detailed survey of approaches, spanning
traditional rule-based systems, classical machine learning, graph neural
networks, spiking neural networks, and autoencoder-based methods. For each
paradigm, we examine the foundational theories, architectural innovations, and
empirical findings with specific attention to performance on the Bank Account
Fraud (BAF) benchmark suite. Finally, we identify five critical research gaps
that this thesis addresses, demonstrating the need for a unified multi-paradigm
comparison under controlled experimental conditions and providing a roadmap for
the contributions that follow in subsequent chapters.

\section{Domain Background}

Financial fraud constitutes one of the most pervasive and economically damaging
categories of cybercrime in the modern digital economy. The scope of financial
fraud is vast, encompassing credit card fraud, identity theft, insurance fraud,
money laundering, and securities manipulation, among other illicit activities.
In the context of credit card fraud alone, global losses exceed \$30 billion
annually, a figure that continues to grow as digital payment systems proliferate
and fraudsters adopt increasingly sophisticated
techniques~\cite{singh2024comprehensive}. Credit card application fraud---the
specific focus of this thesis---occurs when criminals submit applications using
stolen or synthetic identities to obtain credit facilities, which are then
exploited for immediate financial gain before the fraud is detected. The
economic impact of application fraud is particularly severe because each
successful fraudulent application can result in losses ranging from thousands to
tens of thousands of dollars, and organized fraud rings may submit hundreds of
coordinated applications simultaneously~\cite{jesus2022bank}.

The challenge of detecting financial fraud is compounded by several interrelated
factors that distinguish it from conventional classification problems. First,
the extremely low base rate of fraud means that even highly accurate classifiers
can produce an overwhelming number of false positives when deployed at scale,
overwhelming human investigators and eroding confidence in the detection system.
Second, fraud patterns are inherently adversarial: as detection systems evolve,
fraudsters adapt their strategies to circumvent new defenses, creating a
continuous arms race between attackers and defenders that renders static models
obsolete over time. Third, the cost asymmetry between false positives and false
negatives is severe: false positives inconvenience legitimate customers and
erode institutional trust, while false negatives result in direct financial
losses that may be irreversible once the fraudulent activity is discovered.
Fourth, the temporal nature of financial data introduces distributional shift,
as the statistical properties of both legitimate and fraudulent applications
evolve over time due to changes in customer behavior, economic conditions, and
fraudster tactics. These factors collectively make fraud detection one of the
most challenging problems in applied machine learning~\cite{alenzi2020fraud}.

The detection of credit card application fraud presents unique challenges even
within the broader domain of financial fraud detection. Unlike transaction-level
fraud, where each transaction can be evaluated in the context of a known
customer's historical behavior, application fraud involves making decisions
about previously unknown entities with limited historical data. The absence of a
behavioral baseline for new applicants means that detection must rely on the
intrinsic properties of the application itself and its relationships to other
applications and entities. This relational aspect---the fact that fraudulent
applications often share devices, IP addresses, email domains, and other
attributes with other fraudulent applications---is the key insight that
motivates graph-based approaches to fraud detection, which we explore in detail
in Section~\ref{ch:litreview}.

\subsection{The Bank Account Fraud Dataset Suite}

The Bank Account Fraud (BAF) dataset suite~\cite{feedzai2022baf} has emerged as
the primary benchmark for evaluating fraud detection methods under realistic and
challenging conditions. Developed by Feedzai and released in 2022, the BAF suite
was designed to address the limitations of earlier fraud detection benchmarks,
which often used synthetic data, lacked temporal structure, or failed to
represent the extreme class imbalance characteristic of real-world fraud
detection. The base version of the dataset contains approximately one million
application instances, with a fraud prevalence of 1.10\%, yielding a class
imbalance ratio of approximately 89:1. This extreme imbalance is representative
of real-world fraud detection scenarios and poses significant challenges for
standard classification algorithms that assume balanced class distributions. The
BAF dataset is derived from real-world credit card application data that has
been anonymized and preprocessed to protect sensitive information while
preserving the statistical properties relevant to fraud detection research.

A critical feature of the BAF dataset is its temporal structure, which is
essential for realistic evaluation of fraud detection models. The data is
partitioned into training (months 0--4), validation (month 5), and test (months
6--7) splits, reflecting the temporal reality of fraud detection deployments
where models must generalize to future data that may exhibit distributional
shift. This temporal split design stands in sharp contrast to the random
stratified splits commonly used in machine learning research, which tend to
overestimate model performance by allowing information leakage between
temporally adjacent observations~\cite{sun2025objective}. The temporal
separation ensures that models are evaluated on their ability to detect fraud
patterns that may have shifted from the training period, a scenario that is
practically inevitable given the adversarial evolution of fraud strategies and
the natural drift in customer application patterns over time.

The BAF suite comprises six distinct versions, each designed to test a specific
challenge relevant to real-world fraud detection. The \textbf{Base} version
provides the standard evaluation setting with the full feature set and natural
class distribution, serving as the primary benchmark for comparing detection
methods. \textbf{Variant~I} (Group Size Disparity) introduces significant
differences in the number of applications across demographic groups, testing
fairness under population imbalance and challenging models to maintain equitable
performance when some groups are severely underrepresented. \textbf{Variant~II}
(Prevalence Disparity) manipulates the fraud rate across demographic groups,
creating a scenario where the base rate differs substantially across protected
attributes and directly testing the Chouldechova impossibility
result~\cite{chouldechova2017fair}. \textbf{Variant~III} (Prevalence Shift)
introduces a temporal shift in the overall fraud rate between training and test
periods, challenging models to adapt to changing class priors without
retraining. \textbf{Variant~IV} (Temporal Prevalence Shift) combines temporal
drift with group-specific prevalence changes, representing the most realistic
and challenging scenario where both the overall and group-conditional fraud
rates evolve over time. \textbf{Variant~V} (Temporal Separability Shift)
introduces a temporal shift in the separability of legitimate and fraudulent
applications, testing whether models can maintain detection performance when the
feature distributions of the two classes become more similar over time---a
particularly challenging form of distributional shift that reduces the signal
available for discrimination. This comprehensive suite of variants enables
systematic investigation of model behavior under a range of realistic
distributional challenges that are rarely addressed in the fraud detection
literature.

The BAF dataset features comprise 30 attributes, which can be decomposed into 17
continuous or numerical features, 5 categorical features, and 13 binary flags or
indicator variables. The continuous features capture quantitative application
properties such as income, credit history metrics, and application timing, and
require appropriate normalization or scaling before use in neural network
architectures. The categorical features encode discrete attributes such as
employment status and housing type, and must be encoded through techniques such
as one-hot encoding or embedding layers to be processed by deep learning models.
The binary flags indicate the presence or absence of specific risk indicators,
and their sparse, high-dimensional nature presents challenges for models that
assume dense feature representations. This heterogeneous feature space presents
challenges for models that assume homogeneous input distributions, motivating
the exploration of architectures capable of processing mixed data types
effectively. The combination of continuous, categorical, and binary features in
a single dataset reflects the complexity of real-world financial data and makes
the BAF suite a particularly demanding benchmark for deep learning methods that
are typically optimized for homogeneous input types.

\subsection{Fairness Considerations in Fraud Detection}

The deployment of fraud detection systems raises important fairness concerns, as
models may exhibit disparate performance across demographic groups defined by
protected attributes such as age, gender, or ethnicity.
Chouldechova~\cite{chouldechova2017fair} demonstrated the impossibility of
simultaneously achieving equalized odds (equal true positive and false positive
rates across groups) and calibration (predictive parity) when base rates differ
across groups, a condition that is virtually guaranteed in fraud detection due
to the varying prevalence of fraud across demographic segments. This
impossibility result has profound implications: any fraud detection system will
necessarily exhibit some form of unfairness, and the choice of which fairness
metric to optimize involves normative judgments about the relative costs of
different types of errors across groups. The practical consequence is that fraud
detection system designers must make explicit trade-offs between different
fairness criteria, and these trade-offs should be informed by domain-specific
considerations about the relative harms of different error types.

Empirical studies on the BAF dataset have revealed true positive rate (TPR) gaps
of up to 5.5 percentage points across age groups, indicating that certain
demographic segments are systematically less well-served by standard fraud
detection models~\cite{kozodoi2022fairness}. These disparities persist even when
models are trained without explicit access to protected attributes, as proxy
features correlated with group membership can encode discriminatory patterns
that the model learns to exploit. The persistence of fairness disparities in
nominally ``blind'' models underscores the difficulty of achieving equitable
outcomes through data omission alone and motivates the development of fairness-
aware training methods that explicitly constrain model behavior across
demographic groups. Dai et al.~\cite{dai2021fair} have proposed fairness-aware
training objectives for graph-based fraud detection, demonstrating that it is
possible to reduce TPR disparities without substantial degradation in overall
detection performance. However, the interplay between fairness constraints and
the architectural innovations explored in this thesis---including graph neural
networks, spiking neural networks, and autoencoder-based methods---remains
largely unexplored. The tension between detection performance and fairness
underscores the need for multi-metric evaluation frameworks that go beyond
aggregate AUC to assess model behavior across demographic subgroups, a principle
that guides the evaluation methodology of this thesis.

\section{State-of-the-Art in Software Solutions}

Modern financial fraud detection systems increasingly incorporate advanced
analytics alongside traditional tabular processing methods, driven by the
recognition that fraud patterns exploit relational structures that are invisible
to feature-based classifiers. Commercial platforms such as Featurespace,
ThetaRay, and DataVisor have developed proprietary systems that leverage network
analysis and unsupervised anomaly detection for real-time fraud screening,
though the specific algorithms and architectures employed are often not publicly
disclosed. The commercial landscape reflects a growing consensus that single-
paradigm approaches are insufficient for the complexity of modern fraud, and
that effective detection requires the integration of multiple analytical
perspectives. In the academic and open-source domain, several notable frameworks
and benchmarks have emerged that facilitate reproducible fraud detection
research and enable systematic comparison of competing approaches.

The IEEE CIS Fraud Detection competition hosted on Kaggle popularized the use of
rich tabular features for fraud detection, generating a large corpus of
comparative studies on gradient-boosted methods and establishing important
baselines for the field. However, the competition dataset lacked explicit
relational structure, limiting the exploration of graph-based approaches and
focusing research attention on feature engineering and ensemble methods for
tabular data. More recently, graph-based benchmarks such as the Anti-Money
Laundering (AML) dataset and the BAF suite~\cite{feedzai2022baf} have provided
transaction-level data with entity relationships, enabling the systematic
evaluation of graph-based methods. Arulkumaran et
al.~\cite{arulkumaran2022bankguard} proposed BankGuard, a framework integrating
graph analytics with traditional fraud detection pipelines, demonstrating the
practical value of combining relational and feature-based reasoning in
production systems and providing a template for the multi-pipeline architecture
explored in this thesis.

For graph neural network research, PyTorch Geometric (PyG) provides efficient
implementations of various GNN layers including SAGEConv, GATConv, and general
message-passing abstractions that support custom aggregation schemes and
heterogeneous graph processing. The Deep Graph Library (DGL) offers an
alternative framework with native support for heterogeneous graphs, which is
particularly relevant for financial fraud detection where nodes of different
types (applicants, devices, IP addresses) interact through typed edges. For
quantum machine learning, PennyLane and Qiskit provide simulators for
variational quantum circuits that can be integrated with classical neural
network components, enabling research on quantum-inspired architectures without
requiring quantum hardware. For spiking neural networks, frameworks such as
Norse and snnTorch~\cite{eshraghian2021training} provide differentiable Leaky
Integrate-and-Fire (LIF) neuron implementations compatible with PyTorch,
enabling gradient-based training through surrogate gradient methods and
providing the computational infrastructure upon which the SNN pipeline of this
thesis is built. For autoencoder-based anomaly detection, PyTorch and TensorFlow
provide flexible architectures for building VAEs and DAEs, while
XGBoost~\cite{chen2016xgboost} and LightGBM~\cite{ke2017lightgbm} offer
efficient gradient-boosted decision tree implementations that serve as meta-
learners for score-level stacking in hybrid detection pipelines. The
availability of these mature, well-documented frameworks has substantially
lowered the barrier to entry for multi-paradigm fraud detection research and
enables the systematic comparison conducted in this thesis.

\section{Survey of Approaches}

This section presents a detailed survey of approaches to financial fraud
detection, organized by paradigm. We progress from traditional rule-based
systems through classical machine learning to the three deep learning paradigms
central to this thesis: graph neural networks, spiking neural networks, and
autoencoder-based methods. For each paradigm, we review the foundational
methods, recent advances, and specific applications to financial fraud
detection, with emphasis on performance characteristics on the BAF dataset where
available. This progressive organization reflects the historical development of
the field while establishing the context for the multi-paradigm comparison that
is the central contribution of this thesis.

\subsection{Traditional and Rule-Based Approaches}

Traditional fraud detection systems rely on expert-defined rules, heuristic
thresholds, and pattern-matching techniques to identify suspicious applications.
These systems operate by encoding domain knowledge into explicit decision
criteria, such as flagging applications originating from high-risk geographic
regions, those submitted within unusual time windows, or those associated with
known fraudulent devices. Rule-based systems offer several practical advantages
that have sustained their use in the financial industry for decades: they are
fully interpretable, allowing compliance officers and regulators to understand
exactly why an application was flagged; they can be rapidly deployed and
modified in response to emerging fraud patterns without requiring retraining;
and they require no training data, making them immediately applicable in cold-
start scenarios where historical fraud data is
unavailable~\cite{alenzi2020fraud}. In the early days of digital banking, rule-
based systems formed the backbone of fraud detection infrastructure and continue
to serve as a first line of defense in many financial institutions, particularly
for detecting obvious fraud patterns that conform to known templates.

However, rule-based approaches suffer from fundamental limitations that
constrain their effectiveness in modern fraud detection environments. First,
they are inherently reactive, as rules must be crafted in response to previously
observed fraud patterns, leaving them vulnerable to novel attack strategies that
do not match existing rule templates. Sophisticated fraud rings deliberately
vary their tactics to avoid triggering known rules, employing techniques such as
device rotation, geographic distribution, and temporal spacing to evade
detection. Second, the maintenance burden grows combinatorially as the number of
rules increases, leading to rule conflicts, redundancy, and inconsistency that
are difficult to diagnose and resolve. Large financial institutions may maintain
thousands of fraud detection rules, many of which interact in complex and
unpredictable ways. Third, rule-based systems lack the statistical rigor to
handle the probabilistic nature of fraud signals, often producing binary
decisions that fail to account for the uncertainty inherent in fraud assessment.
A rule that flags all applications from a particular IP range, for example,
makes no distinction between high-confidence and low-confidence signals within
that range. Fourth, sophisticated fraudsters can systematically probe and
circumvent rule-based defenses through trial-and-error attacks, as the
deterministic nature of rules makes them predictable once their logic is
inferred. These limitations have motivated the transition toward data-driven
machine learning approaches that can learn complex, non-linear decision
boundaries from historical data while maintaining the flexibility to adapt to
evolving fraud patterns.

Despite their limitations, rule-based systems continue to play a complementary
role in modern fraud detection architectures. Hybrid systems that combine rule-
based pre-screening with machine learning scoring can leverage the strengths of
both approaches: rules provide fast, interpretable filtering for obvious cases,
while machine learning models handle the ambiguous cases that require nuanced
statistical reasoning. This hybrid philosophy extends to the multi-paradigm
approach explored in this thesis, where different models may be deployed in
different stages of a detection pipeline, each contributing its unique strengths
to the overall detection capability. The persistence of rule-based systems in
production environments also serves as a reminder that practical fraud detection
requires not only statistical power but also interpretability, maintainability,
and regulatory compliance---considerations that influence the design choices of
all paradigms examined in this thesis.

\subsection{Classical Machine Learning Approaches}

Classical machine learning methods---including logistic regression, decision
trees, random forests, and gradient-boosted models---have been the workhorses of
fraud detection for over two decades. These methods operate on tabular feature
representations of individual applications, treating each application as an
independent observation and learning decision boundaries that separate
legitimate from fraudulent instances based on feature values. The extensive
literature on classical approaches provides important baselines and insights
that inform the design of more sophisticated deep learning methods, and the
performance of classical models on the BAF dataset establishes the benchmark
against which the deep learning paradigms explored in this thesis must be
measured. The enduring popularity of classical methods in production fraud
detection systems reflects their robustness, interpretability, and strong
performance on tabular data, properties that have been extensively documented in
the machine learning literature~\cite{grinsztajn2022tree}.

\subsubsection{Performance of Classical Models on BAF}

Ahmed and Shamsuddin~\cite{ahmed2022comparative} conducted a comparative study
of classical classifiers on the BAF dataset, reporting accuracy values of
87.91\% for logistic regression, 90.78\% for random forest, and 97.17\% for
decision trees. While these accuracy figures appear impressive at first glance,
they are profoundly misleading due to the extreme class imbalance: a naive
classifier that labels all instances as legitimate would achieve approximately
98.9\% accuracy on the BAF base dataset, rendering accuracy an inappropriate and
deceptive metric for fraud detection evaluation. The high accuracy of the
decision tree in particular reflects overfitting to the majority class rather
than genuine fraud detection capability, as the tree learns to predict the
majority class for most leaf nodes while creating deep, spurious branches that
memorize noise in the minority class. This observation underscores the critical
importance of using class-sensitive metrics such as area under the receiver
operating characteristic curve (AUC), true positive rate at fixed false positive
rate (TPR@FPR), and precision-recall curves when evaluating fraud detection
models, a principle that is consistently applied throughout this thesis.

Islam et al.~\cite{islam2023effects} provided a more nuanced evaluation of
classical methods, reporting that gradient-boosted models (GBM) achieve an AUC
of 0.98 on the BAF dataset, but with a true positive rate below 50\% at
operationally relevant false positive rates. This striking disparity between AUC
and TPR highlights a crucial practical consideration: while the AUC summarizes
detection performance across all possible decision thresholds, operational fraud
detection systems must operate at very low false positive rates (typically 1--
5\%) to avoid overwhelming human investigators with false alerts and degrading
the customer experience for legitimate applicants. At these low FPR operating
points, the detection rate of even the best classical models can be alarmingly
low, leaving a substantial proportion of fraudulent applications undetected.
This finding motivates the exploration of deep learning approaches that may
capture more discriminative patterns, particularly relational patterns that are
invisible to tabular classifiers.

Dong~\cite{dong2025enhanced} further demonstrated the limitations of classical
approaches in a comprehensive evaluation, showing that a random forest model
achieving 98.93\% accuracy on a fraud detection task obtained a fraud recall of
only 0.16\%, meaning that the model detected fewer than 2 in 1,000 fraudulent
applications. This catastrophic failure at the minority class illustrates the
fundamental challenge of class imbalance for classical methods and the
inadequacy of accuracy as an evaluation metric. The results of Dong's study also
highlighted that tree-based ensembles, while generally more robust to imbalance
than single decision trees due to their bagging and boosting mechanisms, still
require explicit imbalance handling strategies to achieve acceptable fraud
detection rates. The consistent pattern across these studies---high accuracy
masking catastrophic minority-class performance---underscores the need for both
improved evaluation methodology and more powerful detection architectures.

\subsubsection{Class Imbalance Handling Strategies}

The extreme class imbalance characteristic of fraud detection datasets has
motivated the development and application of various resampling and loss-
modification strategies, each with distinct mechanisms and trade-offs. The most
widely used resampling methods include SMOTE~\cite{chawla2002smote}, which
generates synthetic minority class examples by interpolating between existing
minority instances and their nearest neighbors in feature space, and
ADASYN~\cite{he2008adasyn}, which adaptively focuses synthetic sample generation
on harder-to-learn minority instances based on their local class distribution.
Loss-based approaches include class weighting~\cite{byrd2019effect}, which
assigns higher misclassification costs to minority class examples to shift the
decision boundary toward better minority-class coverage, and focal
loss~\cite{lin2017focal}, which dynamically down-weights the loss contribution
of well-classified easy examples to focus training on hard, misclassified
instances that are typically concentrated near the decision boundary. Each of
these strategies modifies the effective training distribution in different ways,
and their interaction with different model architectures is an active area of
research.

Uwaoma~\cite{uwaoma2024detecting} conducted a systematic evaluation of five
sampling strategies (SMOTE, ADASYN, random oversampling, random undersampling,
and no sampling) across four model architectures on the BAF dataset, providing
the most comprehensive assessment of imbalance handling strategies for this
benchmark to date. The study revealed that the effectiveness of resampling
strategies is highly architecture-dependent: while resampling can improve the
performance of classical models by providing a more balanced training signal, it
can paradoxically degrade the performance of graph-based methods. Specifically,
ADASYN was found to reduce GraphSAGE AUC from 0.8859 to 0.8748, a
counterintuitive result that highlights the tension between modifying the label
distribution and preserving the relational structure upon which GNNs depend.
When synthetic minority instances are added to the training set, they are
assigned feature vectors that may not correspond to realistic nodes in the
graph, creating artificial neighborhoods that distort the topology and confuse
the message-passing mechanism. This ``resampling paradox'' has significant
implications for the design of multi-paradigm systems, as a preprocessing
strategy that benefits one model type may actively harm another.

Farahat et al.~\cite{farahat2025objective} extended the analysis of imbalance
handling by proposing objective-aware evaluation frameworks that explicitly
account for the operational cost structure of fraud detection, arguing that
standard metrics such as AUC fail to capture the asymmetric costs of false
positives and false negatives in deployment scenarios. Their work motivates the
use of cost-sensitive evaluation metrics such as TPR@5\%FPR, which directly
measures the detection rate at an operationally acceptable false positive rate,
providing a more realistic assessment of model utility. Parkar et
al.~\cite{parkar2024comparative} provided a broader comparative analysis of
preprocessing strategies for fraud detection, confirming that no single
resampling strategy consistently outperforms others across all model
architectures and evaluation metrics, further emphasizing the need for paradigm-
specific tuning and the careful selection of imbalance handling strategies based
on the model architecture and evaluation criteria.

\subsection{Graph Neural Network Approaches}

Graph neural networks (GNNs) have emerged as a powerful paradigm for financial
fraud detection, motivated by the insight that financial fraud is inherently
relational~\cite{pourhabibi2022graph}. Fraud rings share resources---devices, IP
addresses, email domains, and phone numbers---creating network patterns that
become detectable when applicants and their attributes are represented as
interconnected nodes in a graph. This relational perspective transforms the
fraud detection problem from a classification task on isolated feature vectors
to a node classification task on a graph, where the topology itself carries
diagnostic information that is invisible to tabular methods. Recent surveys by
Cheng et al.~\cite{cheng2024graph} and Vanderhulst et
al.~\cite{vanderhulst2024graph} have documented the rapid growth of GNN-based
fraud detection methods, highlighting both the promise and the open challenges
of this approach. The GNN paradigm encompasses a diverse family of
architectures, each with distinct mechanisms for aggregating neighborhood
information and different assumptions about the graph structure, motivating the
systematic comparison conducted in this thesis.

\subsubsection{Graph Construction for Financial Data}

The construction of the graph representation is a critical preprocessing step
that significantly impacts downstream GNN performance and is often overlooked in
comparative studies that focus exclusively on model architecture. For the BAF
dataset, the graph is constructed as a heterogeneous network with approximately
1,000,047 nodes and roughly 14.5 million directed edges spanning 17 attribute
types. Nodes represent both applicants and their associated attributes (devices,
IP addresses, email domains, phone numbers, etc.), while edges encode the
relationships between applicants and their shared attributes. This bipartite-
like structure means that applicants are connected indirectly through shared
attribute nodes, and the resulting graph topology encodes the relational
patterns that are diagnostic of organized fraud. The scale and heterogeneity of
this graph present computational challenges for GNN training, including memory-
efficient neighborhood sampling, mini-batch construction, and handling of
variable-degree nodes.

The design choices involved in graph construction have received increasing
attention in the literature, as researchers have recognized that the graph
representation can be as important as the model architecture in determining
detection performance. Liu et al.~\cite{liu2021heterogeneous} demonstrated that
heterogeneous graph representations, which preserve the type information of
nodes and edges, consistently outperform homogeneous graphs that collapse all
node and edge types into a single category. The type information allows
heterogeneous GNNs to apply type-specific transformations to different edge
categories, enabling the model to distinguish between, for example, a device-
sharing connection and an email-sharing connection, which may carry very
different fraud implications. Pandey et al.~\cite{pandey2022graph} investigated
the impact of feature discretization on graph construction, showing that
continuous features must be binned into discrete categories to serve as
attribute nodes, and that the granularity of discretization affects both graph
density and detection performance. Coarse discretization produces dense graphs
with many shared attribute nodes but limited discriminative power, while fine
discretization produces sparse graphs that may not provide sufficient
connectivity for effective message passing. The interaction between graph
construction strategies and downstream GNN architectures remains an active area
of research, as different architectures may benefit from different graph
representations, and understanding the sensitivity of detection performance to
graph construction choices is essential for the practical deployment of GNN-
based fraud detection systems.

\subsubsection{Spatial GNN Architectures: GraphSAGE}

GraphSAGE~\cite{hamilton2017inductive} introduced the inductive representation
learning paradigm for large-scale graphs, generating node embeddings by sampling
and aggregating features from a node's local neighborhood. Unlike transductive
methods such as DeepWalk and node2vec, which require retraining when new nodes
are added to the graph, GraphSAGE learns an aggregation function that
generalizes to unseen nodes, making it well-suited for the dynamic nature of
financial fraud detection where new applications arrive continuously and must be
classified without retraining the entire model. The original GraphSAGE paper
proposed three aggregation schemes---mean, LSTM, and pooling---each providing
different trade-offs between expressiveness and computational cost. The mean
aggregator computes the element-wise average of neighbor representations,
providing a simple and robust baseline. The LSTM aggregator applies a recurrent
network to a randomly permuted sequence of neighbor features, capturing
sequential dependencies at the cost of order sensitivity. The pooling aggregator
applies a max-pooling operation across neighbor features, capturing the most
salient dimensions of the neighborhood representation.

In the fraud detection domain, GraphSAGE has been applied to capture
neighborhood-level fraud signals through multi-hop information propagation,
allowing the model to detect fraud patterns that extend beyond an individual
applicant's features to include the features of related entities. On the BAF
base dataset, GraphSAGE achieves an AUC of 0.8859 and a TPR@5\%FPR of 52.54\%,
demonstrating that relational information provides substantial discriminative
power beyond what is available from tabular features alone. The mean aggregator,
which computes the element-wise average of neighbor representations, has been
found to be the most effective for fraud detection, likely because the uniform
weighting avoids overfitting to individual neighbor signals in sparse regions of
the graph where the number of neighbors may be small and unrepresentative.

Graph Attention Networks (GAT)~\cite{velivckovic2018graph} extend GraphSAGE by
learning attention weights over neighbor contributions, allowing the model to
assign higher importance to more informative neighbors. GAT-based fraud
detection models can learn to focus on high-risk neighbors (e.g., known
fraudsters sharing a device) while down-weighting less informative connections
(e.g., connections through popular email providers shared by many legitimate
users). However, the attention mechanism introduces additional parameters that
can be prone to overfitting in the sparse, high-dimensional graphs typical of
fraud detection, and the performance gains of GAT over GraphSAGE have been
inconsistent across benchmark datasets. The attention mechanism also adds
computational overhead during both training and inference, which may be
significant for large-scale graphs with millions of edges. These considerations
motivate the use of GraphSAGE as the primary spatial GNN baseline in this
thesis, while acknowledging that attention-based architectures may offer
advantages in specific deployment scenarios where interpretability of neighbor
importance is valued.

\subsubsection{Spectral GNN Architectures: Split GNN with Wavelet Filtering}

Spectral graph theory provides a powerful framework for analyzing graph
structure through the eigendecomposition of the graph Laplacian, enabling the
design of GNN architectures that operate in the frequency domain rather than the
spatial domain. Spectral GNNs apply learnable filters in the frequency domain,
enabling the model to capture both local and global structural patterns through
the spectral properties of the graph. ChebNet~\cite{defferrard2016convolutional}
approximates spectral filters using Chebyshev polynomials of the graph
Laplacian, providing a computationally efficient alternative to exact spectral
filtering that avoids the cost of eigendecomposition. Graph Convolutional
Networks (GCN)~\cite{kipf2017semisupervised} simplify ChebNet to first-order
approximations, resulting in a scalable architecture that aggregates neighbor
features with learned weights but sacrifices the multi-scale filtering
capability of higher-order polynomial approximations. These foundational works
established the viability of spectral filtering for graph representation
learning, but their application to fraud detection has been limited by the
computational cost of spectral decomposition on large-scale graphs and the
difficulty of designing spectral filters that capture fraud-relevant patterns.

The Split GNN architecture~\cite{wu2023splitgnn} represents a significant
advance in spectral GNN design for fraud detection, addressing the limitations
of prior spectral methods through a novel dual-branch filtering mechanism. The
key innovation is the combination of edge classification with dual-branch
spectral filtering using Beta polynomial wavelets~\cite{hammond2011wavelets}.
The edge classifier partitions the graph into subgraphs based on learned edge
semantics, effectively separating same-class connections (homophilic edges,
where connected nodes tend to share the same label) from cross-class connections
(heterophilic edges, where connected nodes tend to have different labels). Each
subgraph is then processed by an independent wavelet filter branch, allowing the
model to apply different spectral transformations to different types of
structural patterns. The Beta wavelet framework enables multi-scale graph signal
decomposition, capturing both fine-grained local patterns and coarse-grained
global structures through a family of polynomial wavelet filters of varying
order. The wavelet order parameter controls the scale of analysis: low-order
wavelets capture local neighborhood patterns, while high-order wavelets capture
broader community structures that may indicate organized fraud rings.

On the BAF base dataset, Split GNN achieves a best AUC of 0.8806, slightly below
GraphSAGE but with different performance characteristics across the precision-
recall trade-off. The dual-branch architecture allows Split GNN to separately
model homophilic connections (links between nodes of the same class, which are
common in shared-device patterns where fraudulent applicants reuse the same
resources) and heterophilic connections (links between nodes of different
classes, which may indicate fraudsters sharing attributes with legitimate users
through popular email providers or widely used IP addresses). This separation is
particularly valuable in fraud detection, where the graph exhibits mixed
homophily: fraudulent applicants are connected to each other through shared
fraudulent resources but also connected to legitimate applicants through common
infrastructure. The wavelet filtering approach also provides natural multi-scale
analysis, with lower-order wavelets capturing local neighborhood structure and
higher-order wavelets capturing broader community patterns that may indicate
organized fraud rings operating across multiple shared resources.

\subsubsection{Quantum-Inspired GNN Architectures}

Quantum machine learning (QML) has emerged as a promising paradigm for capturing
complex patterns through high-dimensional quantum feature spaces, offering the
potential to represent data distributions that are difficult to capture with
classical architectures. Variational quantum circuits (VQCs) parameterized by
trainable rotation angles can encode classical data into quantum states, where
quantum entanglement creates correlations that are difficult to represent
classically and may capture subtle relational patterns in graph-structured data.
Recent work by Herrman et al.~\cite{herrman2022quantum} demonstrated the
application of quantum graph neural networks to classification tasks, showing
that quantum feature maps can capture subtle patterns in graph-structured data
through quantum kernel methods and parameterized quantum circuits. The
theoretical appeal of quantum-enhanced representations lies in the exponential
size of the quantum Hilbert space, which can in principle represent
exponentially many feature combinations with a linear number of qubits.

Innan and Finance~\cite{innan2024quantum} proposed the Quantum Graph Neural
Network (QGNN) architecture, which integrates variational quantum circuits with
classical graph neural network layers for financial fraud detection. The QGNN
uses parameterized quantum circuits to process node features, generating
quantum-enhanced representations that are then propagated through the graph
using classical message-passing operations. This hybrid quantum-classical design
allows the quantum circuits to focus on feature transformation while the
classical GNN handles graph topology, combining the potential advantages of both
paradigms. On the BAF base dataset, QGNN achieves an AUC of 0.8769,
underperforming both GraphSAGE and Split GNN. This result suggests that the
current generation of quantum-inspired architectures, while theoretically
promising, has not yet reached the practical maturity needed to outperform well-
tuned classical GNN methods on large-scale fraud detection tasks. The
underperformance may be attributable to several factors, including the limited
expressiveness of shallow quantum circuits, the difficulty of optimizing
variational parameters in the quantum landscape, and the gap between theoretical
quantum advantage and practical implementation on classical simulators.

A key challenge in QGNN research is the barren plateaus
phenomenon~\cite{mcclean2018barren}, where the gradient of the variational
objective vanishes exponentially with the number of qubits, making optimization
intractable for large-scale circuits. The barren plateaus problem is
particularly severe for the deep, highly parameterized circuits that would be
needed to capture the complex relational patterns in financial fraud graphs, and
it represents a fundamental barrier to scaling quantum-inspired methods to the
problem sizes typical of production fraud detection systems. Current QGNN
implementations address this limitation by using shallow circuits with a small
number of qubits, which limits their expressiveness relative to classical
alternatives but ensures that optimization remains tractable. Our QGNN
implementation uses classical tensor operations to simulate quantum effects,
enabling investigation of quantum-inspired architectures without requiring
quantum hardware. While this simulation approach does not benefit from the
potential quantum advantage of true hardware execution, it provides a controlled
experimental setting for evaluating the architectural contributions of quantum-
inspired designs independently of hardware noise and implementation constraints.
The inclusion of QGNN in the comparison framework of this thesis serves to
establish an empirical baseline for quantum-inspired methods on the BAF
benchmark and to identify the specific conditions under which quantum effects
may or may not provide practical advantages for fraud detection.

\subsection{Spiking Neural Network Approaches}

Spiking neural networks (SNNs) represent the third generation of artificial
neural networks, modeling computation through discrete temporal events (spikes)
rather than continuous activations~\cite{maass1997networks}. Unlike conventional
artificial neurons that communicate real-valued activations through weighted
connections, spiking neurons accumulate incoming synaptic currents into a
membrane potential that evolves over time according to biologically-inspired
differential equations, emitting a discrete spike when the potential exceeds a
threshold. This temporal dynamics enable SNNs to naturally process event-driven
data streams, making them potentially well-suited for transaction-level fraud
detection where temporal ordering and inter-arrival timing carry diagnostic
information. The theoretical appeal of SNNs extends beyond temporal processing
to include energy efficiency on neuromorphic hardware, robustness through
population coding, and the potential for biologically plausible learning rules
that may offer advantages for online adaptation to evolving fraud patterns.

\subsubsection{Mathematical Foundations: The Leaky Integrate-and-Fire Model}

The Leaky Integrate-and-Fire (LIF) neuron is the most widely used spiking neuron
model in the SNN literature, offering a balance between biological realism and
computational tractability that makes it suitable for large-scale simulation and
gradient-based training. The membrane potential dynamics of the LIF neuron are
governed by the following first-order ordinary differential equation, which
describes the subthreshold evolution of the membrane potential as a leaky
integrator of incoming synaptic currents:

\begin{equation}
\tau_m \cdot \frac{dV(t)}{dt} = -\bigl(V(t) - V_{\text{rest}}\bigr) + R \cdot I(t)
\end{equation}

where $V(t)$ denotes the membrane potential at time $t$, $\tau_m$ is the
membrane time constant that controls the rate of potential decay toward the
resting potential, $V_{\text{rest}}$ is the resting potential to which the
membrane decays in the absence of input, $R$ is the membrane resistance that
scales the input current, and $I(t)$ is the input synaptic current at time $t$.
When the membrane potential exceeds a threshold $V_{\text{th}}$, the neuron
emits a spike and the membrane potential is reset to $V_{\text{reset}}$, after
which a refractory period prevents further spiking for a short duration,
mimicking the biological refractory period of real neurons. The discrete-time
approximation of this dynamics, commonly used in simulation frameworks such as
snnTorch, is:

\begin{equation}
V^{(t+1)} = \alpha \cdot V^{(t)} + (1 - \alpha) \cdot \bigl(W \cdot X^{(t)}\bigr)
\end{equation}

where $\alpha = \exp(-\Delta t / \tau_m)$ is the decay factor that determines
how quickly the membrane potential ``forgets'' past inputs, $W$ is the synaptic
weight matrix, and $X^{(t)}$ is the input at time step $t$. The spike generation
is modeled as $S^{(t)} = \Theta(V^{(t)} - V_{\text{th}})$, where $\Theta$ is the
Heaviside step function. The non-differentiability of $\Theta$ presents a
fundamental challenge for gradient-based optimization, as the gradient of the
spike function is zero everywhere except at the threshold, where it is
undefined. This ``dead gradient'' problem is the central computational challenge
of SNN training and has historically limited the depth and complexity of
trainable SNN architectures.

Surrogate gradient methods~\cite{zenke2021remarkable} address this challenge by
replacing the discontinuous spike function with a smooth approximation during
the backward pass, enabling effective gradient propagation while preserving the
discrete spike dynamics during the forward pass. The most commonly used
surrogate gradient is the arc-tangent approximation, which defines the surrogate
derivative as:

\begin{equation}
\frac{\partial S^{(t)}}{\partial V^{(t)}} \approx \frac{1}{\pi} \cdot \frac{1}{1
+ \bigl(\pi \cdot (V^{(t)} - V_{\text{th}})\bigr)^2}
\end{equation}

This approximation provides a smooth, non-vanishing gradient in the vicinity of
the threshold while decaying gracefully for potentials far from the threshold,
preventing gradient explosion and ensuring stable training dynamics. The arc-
tangent surrogate is preferred over alternatives such as the sigmoid or
piecewise linear approximations because it provides a better balance between
gradient magnitude and decay rate, facilitating effective credit assignment
across multiple time steps in deep SNN architectures. Recent advances by
Eshraghian et al.~\cite{eshraghian2021training} have demonstrated that surrogate
gradient training can achieve competitive performance on classification
benchmarks, establishing practical training methodologies for SNNs that open the
door for applications in domains previously dominated by conventional deep
learning, including the financial fraud detection application explored in this
thesis.

\subsubsection{SNNs for Financial Fraud Detection}

The application of SNNs to financial fraud detection is a nascent but rapidly
growing area of research that sits at the intersection of neuromorphic computing
and financial technology. Ahmed~\cite{ahmed2025bioinspired} conducted the first
systematic evaluation of SNNs on the BAF dataset, reporting an AUC of 0.8614 and
a TPR@5\%FPR of 42.21\% on the BAF base version using a multi-layer LIF network
trained with surrogate gradient descent. While these results fall below the
performance of GraphSAGE and classical gradient-boosted methods, the study
demonstrated the viability of SNNs for tabular fraud detection and identified
several architectural and training challenges that limit SNN performance,
providing a foundation for the SNN pipeline developed in this thesis.

Two critical challenges emerged from Ahmed's study that have shaped our
understanding of SNN limitations for fraud detection. First, SNNs exhibit
calibration instability under extreme class imbalance, as the spike rate
distributions of legitimate and fraudulent inputs overlap significantly when the
minority class constitutes only 1.10\% of the training data. The rate coding
scheme used to convert continuous tabular features into spike trains produces
similar firing patterns for both classes, reducing the discriminative capacity
of the spiking representations. The overlap in spike rate distributions means
that the SNN's decision boundary in spike rate space is poorly defined for the
minority class, leading to high false negative rates at the low false positive
rate operating points required by operational fraud detection systems. Second,
SNNs suffer from temporal generalization failure on distributional shift
variants of the BAF dataset, where the temporal drift in feature distributions
between training and test periods disrupts the learned spike timing patterns.
The membrane time constant $\tau_m$, which governs the temporal integration
window of LIF neurons, proves to be a critical hyperparameter: values optimized
for the training period may be inappropriate for the test period, leading to
degraded detection performance. This sensitivity to the membrane time constant
suggests that adaptive or learned temporal dynamics may be necessary for SNNs to
achieve robust performance under distributional shift.

Nasif et al.~\cite{nasif2026reinforcement} proposed a reinforcement-guided
hyper-heuristic approach for SNN optimization, using a reinforcement learning
agent to dynamically select and configure SNN hyperparameters (including
membrane time constants, threshold potentials, and learning rates) during
training. This approach addresses the difficulty of manual hyperparameter tuning
for SNNs, which have a larger and more complex hyperparameter space than
conventional neural networks due to the additional temporal dynamics parameters
of LIF neurons. The reinforcement-guided approach treats hyperparameter
selection as a sequential decision problem, where the agent learns a policy that
maps the current training state to hyperparameter adjustments, enabling adaptive
optimization that responds to the training dynamics. While promising, the
reinforcement-guided approach introduces significant computational overhead due
to the need for repeated policy evaluation and has not yet been evaluated on the
full BAF variant suite, leaving questions about its robustness under
distributional shift and its scalability to large-scale fraud detection
problems.

The potential advantages of SNNs for fraud detection remain largely theoretical
at present, but they motivate continued research into spiking architectures for
financial applications. Rate coding and temporal representations could
potentially capture the sequential nature of transaction streams more naturally
than static feature vectors, encoding the temporal dynamics of application
submission patterns that may be diagnostic of coordinated fraud campaigns. The
event-driven computation paradigm aligns with the irregular arrival patterns of
financial transactions, suggesting potential for low-latency inference in real-
time fraud screening systems where processing speed is critical. Population
coding across neuron ensembles could provide robust representations that are
resilient to adversarial perturbations, an important consideration given the
adversarial nature of fraud and the vulnerability of conventional neural
networks to adversarial examples. Finally, the energy efficiency of SNNs on
neuromorphic hardware may enable edge-deployable fraud screening in scenarios
where computational resources are limited, such as mobile banking applications
or point-of-sale fraud prevention. However, realizing these advantages requires
overcoming the fundamental challenges of training deep SNNs under extreme class
imbalance and distributional shift, challenges that this thesis addresses
through systematic architectural and training methodology exploration.

\subsection{Autoencoder-Based Approaches}

Autoencoders provide a powerful framework for unsupervised anomaly detection by
learning to compress and reconstruct normal data patterns, with reconstruction
errors serving as anomaly scores. The fundamental premise is that an autoencoder
trained predominantly on legitimate transactions will learn to reconstruct
normal patterns accurately while producing high reconstruction errors for
anomalous (fraudulent) transactions that deviate from the learned data manifold.
This unsupervised paradigm is particularly attractive for fraud detection
because it does not require labeled fraud examples during training, addressing
the practical challenge of obtaining reliable fraud labels in operational
settings where fraud confirmation may take weeks or months and the definition of
fraud may evolve over time. The autoencoder framework encompasses several
architectural variants, each with distinct mechanisms for learning the data
manifold and different implications for anomaly detection performance.

\subsubsection{Standard and Variational Autoencoders}

Standard autoencoders consist of an encoder that maps input features to a lower-
dimensional latent representation and a decoder that reconstructs the input from
the latent code. The training objective minimizes the reconstruction error,
typically measured as mean squared error or binary cross-entropy, between the
input and its reconstruction. The encoder and decoder are trained jointly to
minimize this reconstruction loss, with the bottleneck architecture of the
latent space forcing the model to capture the most important aspects of the data
distribution. However, standard autoencoders suffer from an identity mapping
vulnerability: when the model capacity is sufficiently large relative to the
complexity of the input, the autoencoder can learn to reconstruct all inputs---
including anomalies---with low error, effectively memorizing the input rather
than learning the underlying data manifold. This vulnerability is particularly
problematic for fraud detection, where the high dimensionality of financial
features (30 attributes in the BAF dataset) and the relative simplicity of the
fraud signal may allow the autoencoder to reconstruct fraudulent transactions as
effectively as legitimate ones, rendering the reconstruction error uninformative
as an anomaly score.

Variational Autoencoders (VAEs)~\cite{kingma2014auto} address this limitation by
introducing a probabilistic framework that models the data distribution through
a learned latent space. Rather than mapping inputs to single points in the
latent space, VAEs encode inputs as distributions parameterized by mean
$\boldsymbol{\mu}$ and variance $\boldsymbol{\sigma}^2$, from which samples are
drawn using the reparameterization trick that enables gradient-based
optimization through the sampling operation. The training objective maximizes
the evidence lower bound (ELBO), which balances reconstruction quality against a
KL divergence penalty that regularizes the latent space toward the prior
distribution:

\begin{equation}
\mathcal{L}_{\text{ELBO}} = \mathbb{E}_{q(\mathbf{z}|\mathbf{x})}\bigl[\log
p(\mathbf{x}|\mathbf{z})\bigr] - D_{\text{KL}}\bigl(q(\mathbf{z}|\mathbf{x})
\,\|\, p(\mathbf{z})\bigr)
\end{equation}

where the first term encourages faithful reconstruction of the input from the
latent sample, and the KL divergence term regularizes the approximate posterior
$q(\mathbf{z}|\mathbf{x})$ toward the prior $p(\mathbf{z}) =
\mathcal{N}(\mathbf{0}, \mathbf{I})$, preventing the latent space from
overfitting to the training distribution. This regularization is the key
mechanism that mitigates the identity mapping vulnerability: by forcing the
latent space to be smooth and continuous rather than allowing it to memorize
individual instances, the VAE ensures that inputs that deviate from the learned
data manifold are mapped to regions of latent space that produce high
reconstruction errors. The probabilistic nature of VAEs also enables principled
uncertainty quantification, which is valuable for flagging borderline cases for
manual review in operational fraud detection systems.

Obushyni et al.~\cite{obushyni2025vae} conducted a comprehensive evaluation of
VAEs on the BAF dataset, demonstrating that a well-tuned VAE achieves an AUC of
0.8954 on the BAF base version, substantially outperforming a standard
autoencoder that achieves an AUC of only approximately 0.55. The dramatic
performance gap between VAE and standard autoencoder confirms the severity of
the identity mapping problem: the standard autoencoder learns to reconstruct
both legitimate and fraudulent transactions with similar fidelity, rendering the
reconstruction error uninformative as an anomaly score. The VAE's probabilistic
regularization prevents this failure mode by ensuring that the latent space
captures the generative structure of the data rather than memorizing individual
instances, so that fraudulent transactions---which deviate from the learned
generative model---produce higher reconstruction errors and lower log-
likelihoods. Daoud et al.~\cite{daoud2025hybrid} proposed a hybrid VAE-
autoencoder model that combines the probabilistic regularization of VAEs with
the reconstruction fidelity of standard autoencoders, reporting improved
stability across BAF variants and suggesting that the complementary strengths of
probabilistic and deterministic reconstruction can be combined for more robust
anomaly detection. Gouda et al.~\cite{gouda2022deepvae} explored deep VAE
architectures for unsupervised outlier detection, demonstrating that
hierarchical latent spaces with multiple levels of abstraction can capture
multi-scale patterns in financial data, with lower levels capturing fine-grained
feature interactions and higher levels capturing broader distributional
patterns.

\subsubsection{Denoising Autoencoders for Fraud Detection}

Denoising Autoencoders (DAEs)~\cite{vincent2010stacked} take a complementary
approach to the identity mapping problem by learning to reconstruct clean inputs
from corrupted versions. The corruption process forces the autoencoder to learn
robust feature representations that capture the underlying data manifold rather
than memorizing the input, because the model must infer the original values of
corrupted features from the uncorrupted ones, which requires understanding the
inter-feature dependencies that characterize the data distribution. For tabular
fraud data, the most effective corruption strategy is swap noise, where a subset
of input features is randomly replaced with values drawn from the empirical
marginal distribution of the training set. This strategy is particularly
effective for tabular data because it breaks feature-level identity mappings
while preserving the statistical properties of the data, forcing the DAE to
learn inter-feature dependencies that distinguish legitimate from fraudulent
patterns. Anomalous (fraudulent) transactions, which violate the learned inter-
feature relationships, produce high reconstruction errors even without
corruption applied at test time, because the learned dependencies are specific
to the legitimate class distribution.

On the BAF base dataset, a standalone DAE achieves an AUC of only 0.5210, barely
above random guessing and dramatically below the performance of VAE and graph-
based methods. This seemingly poor performance is misleading and requires
careful interpretation: the DAE is learning meaningful representations of the
data manifold, but the reconstruction error alone is insufficient to
discriminate between legitimate and fraudulent transactions when the fraud
signal is subtle relative to the normal variation in the data. The
reconstruction error distribution for fraudulent transactions overlaps heavily
with that for legitimate transactions, making it a poor anomaly score in
isolation. The diagnostic value of the DAE emerges when its learned
representations or reconstruction errors are used as features for a downstream
classifier, rather than as a standalone anomaly detector, a finding that
motivates the hybrid approach explored in this thesis.

Dong~\cite{dong2025enhanced} conducted the most comprehensive evaluation of DAEs
for fraud detection to date, systematically investigating the interaction
between denoising pretraining, resampling strategies, and downstream
classification. A key finding of Dong's study is the resampling paradox: while
ADASYN and SMOTE improve the performance of classical classifiers by providing a
more balanced training signal, they degrade the performance of DAE-based methods
by corrupting the learned data manifold. The synthetic minority instances
generated by resampling methods violate the true data distribution, causing the
DAE to learn a corrupted manifold that does not accurately represent either the
legitimate or fraudulent class. When the DAE is trained on resampled data, the
corruption pattern it learns reflects the artificial distribution rather than
the true data-generating process, reducing the discriminative value of the
reconstruction error. This finding is consistent with the resampling paradox
observed for GNNs and underscores the need for paradigm-specific preprocessing
strategies in multi-model fraud detection systems.

\subsubsection{Stacked Ensemble and Hybrid Approaches}

The complementary strengths of VAE and DAE anomaly detection motivate their
combination through score-level stacking, a meta-learning approach that
leverages the diversity of anomaly signals from multiple autoencoder paradigms.
Unlike feature-level fusion, which concatenates latent representations before
classification, score-level stacking operates on the output anomaly scores,
allowing each autoencoder to optimize its own architecture and training
objective independently. The VAE provides a distribution-aware anomaly score
based on the log-likelihood of the input under the learned generative model,
capturing deviations from the probabilistic data distribution. The DAE provides
a robustness-based anomaly score reflecting the degree to which the input
violates the learned inter-feature dependencies, capturing deviations from the
structural data manifold. These scores capture different aspects of
anomalousness, and their combination through a meta-learner can achieve higher
detection performance than either autoencoder alone by exploiting the orthogonal
information in the two anomaly signals.

Gradient-boosted decision trees---specifically XGBoost~\cite{chen2016xgboost}
and LightGBM~\cite{ke2017lightgbm}---serve as effective meta-learners for score-
level stacking, as they can capture non-linear interactions between anomaly
scores and are robust to the varying scales and distributions of individual
autoencoder outputs. Grinsztajn et al.~\cite{grinsztajn2022tree} demonstrated
that tree-based methods generally outperform deep learning on tabular data,
providing theoretical justification for their use as meta-learners in hybrid
autoencoder-tree pipelines. The tree-based meta-learner can also incorporate
original tabular features alongside autoencoder scores, allowing it to leverage
both the learned representations and the raw feature information for improved
discrimination. Xu et al.~\cite{xu2019neural} proposed neural architecture
search for autoencoder optimization, enabling automated tuning of latent
dimensionality, corruption rate, and architectural choices that are critical for
detection performance.

The hybrid DAE-XGBoost approach achieves an AUC of 0.8752 on the BAF base
dataset, a substantial improvement over the standalone DAE (AUC 0.5210) and
competitive with graph-based methods. This result demonstrates that the DAE
learns useful representations despite its poor standalone performance, and that
the discriminative power of gradient-boosted trees can effectively combine
autoencoder-derived features with original tabular features for improved fraud
detection. The hybrid approach also offers practical advantages in terms of
interpretability, as the feature importance scores of the gradient-boosted meta-
learner can reveal which autoencoder scores and original features contribute
most to the detection decision. Pombal et al.~\cite{pombal2022understanding}
provided a systematic study of ensemble strategies for anomaly detection,
confirming that score-level stacking consistently outperforms both individual
models and feature-level fusion approaches, particularly under the extreme class
imbalance conditions characteristic of fraud detection.

\subsection{Comparative Analysis of Approaches}

The survey of approaches presented in the preceding sections reveals a complex
landscape in which no single paradigm consistently dominates across all
evaluation metrics and dataset variants. Table~\ref{tab:litreview-comparison}
summarizes the representative performance of each approach on the BAF base
dataset, highlighting the trade-offs between AUC, TPR@5\%FPR, and architectural
complexity that practitioners must navigate when selecting detection methods.

\begin{table}[htbp]
\centering
\caption{Representative performance of fraud detection approaches on the BAF
base dataset. Entries marked ``---'' indicate metrics not reported in the
original study.}
\label{tab:litreview-comparison}
\begin{tabular}{lcc}
\toprule
\textbf{Approach} & \textbf{AUC} & \textbf{TPR@5\%FPR} \\
\midrule
Gradient-Boosted Trees (GBM) & 0.98 & $<$50\% \\
VAE & 0.8954 & --- \\
GraphSAGE & 0.8859 & 52.54\% \\
Split GNN & 0.8806 & --- \\
QGNN & 0.8769 & --- \\
DAE-XGBoost (Hybrid) & 0.8752 & --- \\
SNN (LIF) & 0.8614 & 42.21\% \\
Standalone DAE & 0.5210 & --- \\
Standard Autoencoder & $\sim$0.55 & --- \\
\bottomrule
\end{tabular}
\end{table}

Several important observations emerge from this comparative analysis that inform
the design of this thesis. First, classical gradient-boosted methods achieve the
highest AUC overall (0.98), but their TPR at operationally relevant FPR levels
is below 50\%, indicating that a significant proportion of fraudulent
applications evade detection in practical deployment scenarios where the false
positive budget is tightly constrained. This paradox---high AUC with low
operational TPR---reflects the fact that AUC averages performance across all
possible thresholds, including many that are operationally irrelevant, and
motivates the use of TPR@5\%FPR as a complementary metric throughout this
thesis.

Second, GraphSAGE achieves the highest TPR@5\%FPR (52.54\%) among the deep
learning methods compared, demonstrating the value of relational information for
detecting the most subtle fraud patterns that escape tabular classifiers. The
superior TPR at low FPR suggests that the graph topology encodes information
that is particularly valuable for identifying the borderline cases that are most
challenging for tabular methods, and that the message-passing mechanism of
spatial GNNs effectively propagates this information to the nodes that need it
most.

Third, the VAE achieves a surprisingly high AUC (0.8954) as a standalone
unsupervised method, outperforming several supervised approaches and indicating
that the probabilistic latent space captures fraud-relevant structure even
without explicit supervision. This result suggests that unsupervised methods may
have a valuable role to play in multi-paradigm fraud detection systems, either
as standalone detectors or as components of hybrid ensembles that combine
supervised and unsupervised signals.

Fourth, the dramatic gap between standalone DAE (AUC 0.5210) and hybrid DAE-
XGBoost (AUC 0.8752) underscores the importance of downstream classification for
unlocking the discriminative value of autoencoder representations and
demonstrates that the DAE's failure as a standalone detector does not reflect a
failure of representation learning but rather a failure of the reconstruction
error as a direct anomaly score for tabular fraud data.

Fifth, SNNs currently lag behind other deep learning paradigms, with the lowest
AUC (0.8614) and TPR@5\%FPR (42.21\%) among the compared methods, indicating
that significant architectural and training innovations are needed to make SNNs
competitive for fraud detection. The calibration instability and temporal
generalization failure identified by Ahmed~\cite{ahmed2025bioinspired} represent
specific technical challenges that this thesis addresses through systematic
experimentation with SNN architectures and training methodologies.

Finally, the incompleteness of the comparison table---with several TPR@5\%FPR
values unreported---reflects a broader problem in the literature: the absence of
a unified evaluation protocol that reports a consistent set of metrics across
all paradigms. This gap motivates the systematic multi-paradigm comparison
conducted in this thesis, which evaluates all approaches under identical
conditions using a comprehensive suite of metrics.

\section{Identification of Research Gaps}

The comprehensive literature review presented in this chapter reveals several
critical gaps in the current state of research on deep learning for financial
fraud detection. Each gap represents an opportunity for methodological
innovation and empirical investigation that this thesis addresses. In this
section, we identify and discuss five principal research gaps, followed by a
summary of how this thesis fills these voids and advances the state of the art.

\subsection{Gap 1: Absence of Multi-Paradigm Comparison Under Unified Conditions}

The existing literature on deep learning for fraud detection is fragmented
across individual paradigms, with studies typically evaluating a single model
family in isolation using dataset-specific preprocessing and evaluation
protocols. Pourhabibi~\cite{pourhabibi2022graph} and Cheng et
al.~\cite{cheng2024graph} have both highlighted the lack of systematic cross-
paradigm comparisons in their surveys of GNN-based fraud detection, noting that
the absence of controlled comparisons makes it impossible to determine whether
observed performance differences reflect genuine architectural advantages or
merely differences in hyperparameter tuning, preprocessing pipelines, and
evaluation methodology. This fragmentation is particularly problematic for
practitioners who must select among competing paradigms for deployment, as the
relative strengths and weaknesses of GNN, SNN, and autoencoder-based approaches
have never been characterized under identical experimental conditions.

The absence of unified comparison also obscures the complementary potential of
different paradigms. As demonstrated in our comparative analysis, GraphSAGE
excels at TPR@5\%FPR, VAEs achieve high AUC through probabilistic modeling, and
hybrid DAE-XGBoost systems leverage autoencoder representations for
discriminative classification. A multi-paradigm system that combines these
complementary strengths could potentially achieve performance that exceeds any
single paradigm, but the design of such a system requires a detailed
understanding of how the paradigms differ in their error patterns, failure
modes, and sensitivity to preprocessing choices. The lack of unified comparison
also means that the resampling paradox---where ADASYN helps classical models but
hurts GNNs and DAEs---cannot be assessed in the context of a complete multi-
paradigm system, limiting the practical guidance available for system design.
This thesis provides the first systematic multi-paradigm comparison under
unified experimental conditions on the BAF dataset suite, establishing the
empirical foundation for principled paradigm selection and combination.

\subsection{Gap 2: Insufficient Exploitation of Relational Structure in Anomaly
Detection}

Autoencoder-based anomaly detection methods, including both VAEs and DAEs,
operate exclusively on tabular feature representations and fail to exploit the
relational structure that graph-based methods leverage for fraud detection. This
represents a significant missed opportunity, as the graph topology encodes
relational patterns---such as shared devices, common IP addresses, and
overlapping contact information---that are among the most reliable indicators of
organized fraud. The current separation between graph-based methods (which
exploit relational structure but may not capture fine-grained feature
interactions) and autoencoder-based methods (which capture feature interactions
but ignore relational structure) suggests that a combined approach could achieve
superior performance by leveraging both sources of information simultaneously.

The integration of relational information into autoencoder architectures is not
straightforward and presents several design challenges. Naive approaches such as
concatenating graph-derived features (e.g., GraphSAGE embeddings) with original
tabular features before autoencoder input may fail because the autoencoder can
learn to reconstruct the graph features independently of the tabular features,
diluting the relational signal and reducing the complementary value of the
combined representation. More sophisticated approaches, such as graph-
regularized autoencoders that enforce consistency between the autoencoder latent
space and the graph topology, or conditional VAEs that generate reconstructions
conditioned on neighborhood information, have been proposed in the broader
representation learning literature but have not been systematically evaluated
for financial fraud detection. This thesis investigates the potential for
relational information to improve autoencoder-based anomaly detection through
score-level combination with graph-based methods, providing the first empirical
assessment of cross-paradigm complementarity on the BAF dataset and establishing
whether the relational and feature-based signals are indeed orthogonal and
combinable.

\subsection{Gap 3: Limited Understanding of Denoising Autoencoders for Tabular
Fraud Data}

While denoising autoencoders have a well-established theoretical foundation in
the image domain~\cite{vincent2010stacked}, their application to tabular fraud
data remains poorly understood. The swap noise corruption strategy that is
effective for tabular data introduces unique challenges that differ
fundamentally from the Gaussian noise or masking corruption used in image-based
DAEs. In the image domain, spatial locality ensures that corrupted pixels can be
inferred from their neighbors, but tabular features lack spatial structure, and
the inter-feature dependencies that enable denoising are learned purely from
statistical correlations in the data. The optimal corruption rate, the
interaction between corruption rate and class imbalance, and the relationship
between DAE latent dimensionality and detection performance have not been
systematically characterized for fraud detection. Dong~\cite{dong2025enhanced}
provided the most comprehensive evaluation to date, but several questions remain
open, including the sensitivity of DAE performance to the specific features
selected for corruption, the impact of corruption rate scheduling during
training, and the effectiveness of DAE-derived features for detecting specific
fraud subtypes.

Furthermore, the stark contrast between standalone DAE performance (AUC 0.5210)
and hybrid DAE-XGBoost performance (AUC 0.8752) raises fundamental questions
about what the DAE actually learns and why its representations are useful only
when combined with a downstream classifier. Understanding the nature of DAE
representations for fraud data---whether they capture class-relevant structure,
feature interactions, or data manifold geometry---is essential for designing
more effective autoencoder architectures and for determining the conditions
under which DAE-based methods are most likely to succeed. The identity mapping
vulnerability, while mitigated by the denoising objective, may persist in
attenuated form, and the conditions under which the DAE fails to learn
discriminative representations remain unclear. This thesis contributes to this
understanding through systematic ablation studies that isolate the contributions
of different DAE design choices to downstream detection performance, providing
practical guidance for the deployment of DAE-based methods in fraud detection
systems.

\subsection{Gap 4: Temporal Robustness Across Distributional Shifts}

The temporal evaluation framework of the BAF dataset suite, with its six
variants designed to test different aspects of distributional shift, reveals a
critical gap in the current literature: the temporal robustness of fraud
detection models is insufficiently characterized. Sun et
al.~\cite{sun2025objective} demonstrated that models optimized for performance
on the BAF base version often exhibit severe performance degradation on the
temporal shift variants (particularly Variant~IV and Variant~V), yet most
published studies report results only on the base version or on random
stratified splits that do not reflect the temporal reality of fraud detection.
This practice leads to an overly optimistic assessment of model capability and
fails to identify the architectural and training choices that promote temporal
robustness, creating a misleading picture of model readiness for production
deployment.

The challenge of temporal robustness is paradigm-specific, with each paradigm
exhibiting distinct failure modes under distributional shift that reflect its
architectural assumptions and learning mechanisms. For GNNs, distributional
shift may alter the graph topology (as new types of attribute nodes emerge over
time) and the feature distributions of existing nodes, requiring inductive
generalization beyond the training graph structure and potentially disrupting
the message-passing patterns that the model has learned. For SNNs, temporal
shift disrupts the learned spike timing patterns and membrane dynamics, as the
optimal membrane time constant may change with the evolving data distribution,
causing the temporal integration window to become misaligned with the true
temporal structure of the data. For autoencoders, distributional shift may cause
the learned data manifold to diverge from the test distribution, leading to
systematically biased reconstruction errors that affect the calibration of
anomaly scores. Understanding these paradigm-specific failure modes and
developing mitigation strategies---such as temporal regularization, domain
adaptation, or continuous learning---is essential for deploying fraud detection
models in production environments where the data distribution is non-stationary.
This thesis evaluates all three paradigms across the full BAF variant suite,
providing the first comprehensive characterization of temporal robustness across
deep learning paradigms for fraud detection.

\subsection{Gap 5: Integration of Software Engineering Practices with Multi-
Model Fraud Detection}

The deployment of multi-paradigm fraud detection systems in production
environments raises software engineering challenges that have received virtually
no attention in the academic literature. While individual model architectures
have been extensively studied, the engineering of systems that combine multiple
paradigms---including model versioning, reproducible training pipelines,
consistent evaluation protocols, and maintainable codebases---has been largely
neglected. The practical challenges of deploying and maintaining a multi-model
system in a financial institution include managing dependencies between model
components, ensuring reproducibility across hardware and software environments,
implementing robust monitoring and alerting for model degradation, and
maintaining audit trails for regulatory compliance. These engineering
considerations are not merely incidental but are essential for translating
academic research into operational systems that deliver reliable fraud detection
at scale.

The gap between academic model development and production system engineering is
particularly acute for multi-paradigm systems, where each paradigm may require
different preprocessing, training infrastructure, and inference pipelines. GNNs
require graph construction and neighborhood sampling, which involve complex data
structures and sampling algorithms that must be carefully validated for
correctness and efficiency. SNNs require temporal simulation and surrogate
gradient computation, which introduce additional computational steps and
hyperparameters that must be managed and tuned. Autoencoders require corruption
scheduling and latent space regularization, which add complexity to the training
pipeline and may interact with other preprocessing steps. Integrating these
diverse requirements into a cohesive, maintainable system demands careful
software architecture design and adherence to engineering best practices such as
modular design, comprehensive testing, and configuration management. This thesis
addresses this gap by developing a well-engineered multi-pipeline system that
demonstrates the feasibility of integrating GNN, SNN, and autoencoder paradigms
within a unified software framework, providing a template for the practical
deployment of multi-model fraud detection systems that can be maintained,
extended, and audited by engineering teams.

\subsection{Summary: How This Thesis Fills the Void}

The five research gaps identified above collectively define the contribution
space of this thesis and establish the motivation for the multi-paradigm
investigation that follows. Gap~1 motivates the need for a unified multi-
paradigm comparison, which this thesis provides through systematic evaluation of
GraphSAGE, Split GNN, QGNN, LIF-based SNN, VAE, and DAE-XGBoost pipelines on the
complete BAF variant suite under identical experimental conditions, including
consistent preprocessing, evaluation metrics, and hyperparameter tuning
procedures. Gap~2 motivates the exploration of cross-paradigm complementarity,
which this thesis investigates through score-level combination of graph-based
and autoencoder-based anomaly scores, testing the hypothesis that relational and
feature-based signals provide orthogonal information that can be combined for
improved detection. Gap~3 motivates the systematic study of DAE design choices
for tabular fraud data, which this thesis addresses through comprehensive
ablation experiments that characterize the sensitivity of DAE performance to
corruption rate, latent dimensionality, and architectural configuration. Gap~4
motivates the evaluation of temporal robustness across distributional shifts,
which this thesis provides through evaluation on all six BAF variants with
temporal train-test splits, characterizing the paradigm-specific failure modes
and identifying architectural choices that promote robustness. Gap~5 motivates
the integration of software engineering practices, which this thesis
demonstrates through the design and implementation of a modular, reproducible
multi-pipeline system that serves as both a research platform and a template for
production deployment.

By addressing these gaps in a single, coherent body of work, this thesis makes
both methodological and practical contributions to the field of financial fraud
detection. The methodological contributions include the first systematic multi-
paradigm comparison, the characterization of cross-paradigm complementarity, and
the identification of paradigm-specific failure modes under distributional
shift. The practical contributions include design recommendations for multi-
paradigm fraud detection systems, engineering patterns for multi-model
deployment, and empirical guidelines for paradigm selection and combination in
different operational contexts. Together, these contributions advance the state
of the art from isolated single-paradigm studies toward an integrated, multi-
paradigm understanding of deep learning for financial fraud detection, providing
both the empirical evidence and the engineering infrastructure needed to
translate research advances into operational fraud detection systems.
